% =====================================================================
%  Prefacio — Econometría Financiera Aplicada (Jeel Cueva)
% =====================================================================
\chapter*{Prefacio}
\addcontentsline{toc}{chapter}{Prefacio}
\markboth{Prefacio}{Prefacio}

Este libro nació de una necesidad concreta: un texto de econometría financiera en español que llevara al
lector desde las nociones básicas de precios y rendimientos hasta los modelos que hoy se usan en la gestión
de riesgos, la valoración de activos y el pronóstico, sin renunciar al rigor de las demostraciones ni al
trabajo con datos reales. Cada resultado teórico se demuestra paso a paso y cada técnica se aplica a series
reales del Perú y del mundo: el índice de la Bolsa de Valores de Lima y el tipo de cambio del Banco Central
de Reserva del Perú, acciones peruanas que cotizan en Nueva York, el petróleo, el S\&P 500, las carteras de
Fama y French y la curva de rendimientos de los Estados Unidos, entre otras.

El libro tiene veintiún capítulos y tres apéndices. Los primeros capítulos tratan los rendimientos, sus
hechos estilizados y la regresión en finanzas; les siguen los modelos de series de tiempo, la volatilidad
condicional, los modelos multivariados, el valor en riesgo, los modelos de factores, las tasas de interés,
los estudios de eventos, los cambios de régimen, las cópulas, los datos de panel y el aprendizaje
automático. El último capítulo integra todo en un proyecto completo. Cada capítulo incluye ejercicios
resueltos y ejercicios propuestos.

Todos los ejercicios resueltos siguen la misma estructura: el enunciado, el planteamiento y los pasos de
la solución, la salida real del programa en R, una tabla \emph{Cómo leer la salida} que interpreta cada
indicador que aparece en ella y una conclusión. La tabla está pensada para que el lector no solo obtenga
un número, sino que sepa qué significa, cuándo es grande o pequeño y qué decisión permite tomar.

Para que el libro no se hiciera excesivamente extenso, los programas no se imprimen. El código completo en
R y en Python de cada ejercicio, la salida que produce en ambos lenguajes y todas las bases de datos están
en el repositorio del libro en GitHub:
\begin{center}
\begin{tabular}{@{}c@{\qquad}l@{}}
\raisebox{-0.5\height}{\qrcode[height=2.6cm]{\sitioweb}} &
\begin{minipage}{0.62\linewidth}\raggedright
\href{\sitioweb}{\nolinkurl{\sitioweb}}\\[6pt]
\small Escanee el código con el teléfono o haga clic en la dirección. Antes de la salida de cada ejercicio
hay enlaces directos a sus programas en el repositorio.
\end{minipage}
\end{tabular}
\end{center}
El repositorio está organizado como el proyecto del libro. La carpeta \texttt{codigo} contiene un programa
por ejercicio en R (\texttt{codigo/R}) y en Python (\texttt{codigo/Python}); por ejemplo, los del
Ejercicio~1.1 son \texttt{cap01\_ej01.R} y \texttt{cap01\_ej01.py}. La carpeta \texttt{salidas} contiene la
salida de cada programa en ambos lenguajes, y la carpeta \texttt{datos} reúne los archivos CSV con un
archivo \texttt{FUENTES.txt} que documenta sus fuentes y fechas de descarga. Todo el material puede
descargarse en un solo archivo con el botón \emph{Code} $\rightarrow$ \emph{Download ZIP} de GitHub.
Las cifras del libro son exactamente las que producen esos programas, de
modo que el lector puede reproducir cualquier resultado y modificar los programas para aplicarlos a sus
propios datos.

Los datos son reales y públicos. Las series peruanas provienen de BCRPData y de Yahoo Finance; las
internacionales, de bases académicas y oficiales que se citan en cada capítulo. Al trabajar con ellas
aparecieron errores en las propias fuentes, como un dato erróneo del índice bursátil o un dividendo
registrado dos veces, que el libro no oculta: se detectan, se explican y se corrigen, porque reconocerlos
es parte del oficio del econometrista.

El libro está pensado tanto para cursos universitarios de pregrado avanzado y maestría como para el
estudio personal. Se supone que el lector conoce la estadística y el álgebra lineal de un primer curso;
los resultados que se usan sin demostrar en el texto principal se reúnen en los Apéndices~A y~B.

Agradeceré a los lectores que me hagan llegar sus comentarios y los errores que encuentren, para
corregirlos en futuras ediciones.

\bigskip
\noindent Lima, noviembre de 2026\hfill Jeel Cueva
